\chapter{Pay Levels by Role, Firm Type and Seniority}\label{ch:in:pay-levels-by-role-firm-type-and-seniority}

Every public source of pay sees one slice. US visa filings see the base salaries offered for sponsored positions, and
nothing of bonuses. The government's occupational survey sees wages by industry and occupation, but excludes most
bonuses and caps what it prints. Filed accounts see total staff cost per head, bonuses included, but average away the
roles. Bank remuneration reports count only the people paid more than a million euros. Put side by side, the slices
disagree by a factor of ten and more, and each is right about what it measures: in 2025 the median \omterm{def:in:how-pay-works:base}{base salary} offered
in sponsored applications at the market makers of this book was \$175\,000, while one of them spent \$2.1 million per
UK employee on staff. This chapter builds ranges from each slice with its caveat, by role, by kind of employer and by
seniority, and then joins them into the one number a reader most wants and no source prints: how total pay relates to
base.

\section{What each source measures}

Five public sources carry pay figures for the industry. Each measures a different quantity over a different population,
and none measures what a reader means by ``pay''.

\begin{definition}[Labor condition application, prevailing wage]\label{def:in:pay-levels-by-role-firm-type-and-seniority:lca}
A \emph{labor condition application}\index{labor condition application} (LCA) is the attestation a US employer files
with the Department of Labor before sponsoring a worker on an H-1B, H-1B1 or E-3 visa: the job title, occupation code,
worksite, the offered wage or wage range, and the prevailing wage for the occupation and area. The \emph{prevailing
wage}\index{prevailing wage} is the wage the Department's rules set as the minimum for the occupation, area and level;
the employer must pay the greater of it and the actual wage it pays similar employees.
\end{definition}

\begin{itemize}
\item \textbf{Labour condition applications} give the offered annual base of each sponsored position. A bonus counts
toward the required wage only ``if their payment is assured'', not when it is ``contingent on some event such as the
employer's annual profits'', so the offered wage is \omterm{def:in:how-pay-works:base}{base salary}. The population is sponsored positions, not all staff.
\item \textbf{The occupational survey} (OEWS) gives wage percentiles by industry and occupation. Its wages include
``incentive pay, including commissions and production bonuses'' but exclude ``nonproduction bonuses'', which is where a
discretionary bonus falls, and it marks the highest cells as top-coded.
\item \textbf{Filed accounts} give \omterm{def:in:reading-a-trading-firms-accounts:staff}{staff costs} over \omterm{def:in:reading-a-trading-firms-accounts:staff}{average headcount} (chapter~11): every employee, every element of
pay the year's accounts charge, and the employer's social costs, for one entity.
\item \textbf{Pay-ratio disclosures} of US listed firms give the total pay of their median employee.
\item \textbf{Remuneration disclosures} of banks and investment firms give counts and totals for the best-paid:
the Pillar~3 tables of chapter~11 and the European Banking Authority's high-earner reports.
\end{itemize}

\begin{definition}[Top-coding, small-cell suppression]\label{def:in:pay-levels-by-role-firm-type-and-seniority:suppress}
\emph{Top-coding}\index{top-coding} is the replacement of values above a threshold by a marker, so that a survey's
highest cells show only that they exceed it. \emph{Small-cell suppression}\index{small-cell suppression} is the rule of
not publishing a statistic computed from too few observations, because it would describe identifiable individuals or
be too noisy to use.
\end{definition}

The occupational survey's field descriptions still define \# as ``a wage equal to or greater than \$115.00 per hour or
\$239\,200 per year'', although the May 2025 file prints values above that line (none of the cells this book uses is
marked); this chapter applies its own suppression to the visa filings: no cell is printed from fewer than ten applications, or from
fewer than three employers, because a cell drawn from one or two employers would publish a firm's pay.

\begin{definition}[Median-employee pay, high earner]\label{def:in:pay-levels-by-role-firm-type-and-seniority:median}
\emph{Median-employee pay}\index{median-employee pay} is the annual \omterm{def:in:how-pay-works:base}{total compensation} of the median employee of a US
listed company, which its proxy statement discloses with the chief executive's pay ratio. A \emph{high
earner}\index{high earner}, in EU banking statistics, is an individual paid one million euros or more in a financial
year; banks and investment firms report their numbers and pay to their supervisors, and the European Banking Authority
publishes the aggregates.
\end{definition}

\section{Building ranges from labour-condition filings}

The Department of Labor publishes every application it decides, a quarter at a time, in files of up to a few hundred
thousand rows. The chapter reads the four quarterly files of fiscal years 2021 and 2025, deduplicated by case number.

\begin{definition}[Standard Occupational Classification, wage level]\label{def:in:pay-levels-by-role-firm-type-and-seniority:soc}
The \emph{Standard Occupational Classification}\index{Standard Occupational Classification} (SOC) is the US federal
statistical standard that assigns every worker to one of 867 detailed occupations; 13-2099.01 is ``Financial
Quantitative Analysts''. A \emph{wage level}\index{wage level} is one of four prevailing-wage levels an application is
assigned to: I (entry), II (qualified), III (experienced) and IV (fully competent), chosen from the job's requirements.
\end{definition}

\begin{method}[Ranges from labour condition applications]\label{met:in:pay-levels-by-role-firm-type-and-seniority:lca}
\begin{enumerate}
\item Stream each file row by row, reading only the needed columns; never the contact fields.
\item Keep certified, full-time applications.
\item Classify the employer by name with rules whose business model is sourced in an earlier chapter (twelve market
makers, seven systematic funds, four platforms, nine banks, one asset manager, four exchanges), and the role by the job
title, or by the SOC code when the title is generic.
\item Annualise the lower end of the offered wage by its unit (year, month, two weeks, week, hour).
\item Group by fiscal year, kind of employer, role and \omterm{def:in:pay-levels-by-role-firm-type-and-seniority:soc}{wage level}; compute percentiles and a bootstrap interval for the
median; suppress cells under ten applications or three employers.
\end{enumerate}
\end{method}

\begin{dated}{2025-09}{What the visa filings contain for this book's employers}\label{dat:in:pay-levels-by-role-firm-type-and-seniority:counts}
\footnotesize
\begin{tabular}{@{}lrrrr@{}}
\toprule
kind of employer & applications FY2021 & employers & applications FY2025 & employers\\
\midrule
bank & 3\,539 & 9 & 7\,298 & 9\\
market maker & 259 & 10 & 367 & 10\\
systematic fund & 143 & 5 & 210 & 7\\
multi-manager platform & 53 & 3 & 102 & 3\\
exchange & 228 & 4 & 462 & 4\\
asset manager & 200 & 1 & 369 & 1\\
\bottomrule
\end{tabular}

\smallskip\noindent Certified full-time applications with a role in the chapter's families, US fiscal years (October to
September), from the Department of Labor's disclosure files. The asset manager's cells come from one employer and are
all suppressed.
\end{dated}

The banks file twenty times as many applications as the market makers, because they employ many more people in
technology; a statistic pooled over the industry would be a bank statistic. Every range below is therefore by kind of
employer.

\section{Ranges by role, firm type and seniority}

\begin{center}\footnotesize
\begin{tabular}{@{}lrrrrr@{}}
\toprule
median offered base, FY2025 (\$ thousand) & systematic fund & platform & market maker & bank & exchange\\
\midrule
all roles & 215.0 & 190.5 & 175.0 & 154.3 & 130.8\\
quantitative researcher or analyst & 220.0 & 190.0 & 175.0 & 158.1 & 107.1\\
software engineer & 205.5 & 190.0 & 175.0 & 155.7 & 140.4\\
trader & -- & -- & 195.0 & 235.0 & --\\
\bottomrule
\end{tabular}

\smallskip Cells marked -- are suppressed. Data: \texttt{data/industry/lca\_ranges.csv}.
\end{center}

The order by kind of employer is the same for both main roles: systematic funds offer the highest base, then the
platforms, the market makers, the banks and the exchanges, with about \$84\,000 between the top and the bottom for all
roles together (\cref{fig:in:pay-levels-by-role-firm-type-and-seniority:kind}). Within a kind, the spread is wide: the
market makers' interquartile range for quantitative researchers runs from \$132\,400 to \$237\,500. From fiscal 2021
to 2025 the median across all roles rose 22.9\% at the systematic funds and at the banks, 18.7\% at the platforms, 16.7\%
at the market makers and 7.9\% at the exchanges; US consumer prices rose 18.8\% over the same calendar years, so in real
terms offered base was flat to slightly up, and down at the exchanges.

\begin{omfigure}
\begin{tikzpicture}
\begin{axis}[width=10.5cm,height=5.8cm,xmin=0.4,xmax=5.6,ymin=60,ymax=320,xtick=data,
  xticklabels from table={figdata/industry/14-pay-levels-by-role-firm-type-and-seniority/base_by_kind.csv}{kind},
  xticklabel style={font=\footnotesize,text width=1.8cm,align=center},yticklabel style={font=\footnotesize},
  ylabel={\small offered base, \$ thousand},
  grid=major,grid style={gray!20},table/col sep=comma,unbounded coords=jump,
  legend style={legend cell align=left,font=\scriptsize,at={(0.5,-0.3)},anchor=north,/tikz/every even column/.append style={column sep=8pt}},
  legend columns=3]
\addplot+[only marks,mark=*,color=omDef,error bars/.cd,y dir=both,y explicit]
  table[x expr=\thisrow{pos}-0.18,y=all,y error plus=all_up,y error minus=all_dn]{figdata/industry/14-pay-levels-by-role-firm-type-and-seniority/base_by_kind.csv};
\addplot+[only marks,mark=square*,color=omThm,error bars/.cd,y dir=both,y explicit]
  table[x=pos,y=qr,y error plus=qr_up,y error minus=qr_dn]{figdata/industry/14-pay-levels-by-role-firm-type-and-seniority/base_by_kind.csv};
\addplot+[only marks,mark=triangle*,color=black!70,error bars/.cd,y dir=both,y explicit]
  table[x expr=\thisrow{pos}+0.18,y=se,y error plus=se_up,y error minus=se_dn]{figdata/industry/14-pay-levels-by-role-firm-type-and-seniority/base_by_kind.csv};
\legend{all roles,quantitative researcher or analyst,software engineer}
\end{axis}
\end{tikzpicture}

\omcaption{Median offered annual base in US labour condition applications of fiscal 2025, with the interquartile range,
by kind of employer and role. Base only: bonuses are not in these filings. Data: \texttt{data/industry/lca\_ranges.csv},
from the Department of Labor's disclosure files through \texttt{in\_lca\_derive.py}.}\label{fig:in:pay-levels-by-role-firm-type-and-seniority:kind}
\end{omfigure}

The \omterm{def:in:pay-levels-by-role-firm-type-and-seniority:soc}{wage level} is the only seniority measure the filings carry, and it is a weak one where pay is far above the floor
(\cref{fig:in:pay-levels-by-role-firm-type-and-seniority:level}). At the banks and the exchanges the median rises with
the level, as the Department's definitions intend: at the banks from \$100\,000 at level~I to \$169\,900 at level~IV.
At the market makers the median is \$84\,800 at level~I and \$180\,000 to \$185\,000 at levels II to IV: the employer
picks the level that describes the job's requirements, and when the offer is far above every level's \omterm{def:in:pay-levels-by-role-firm-type-and-seniority:lca}{prevailing wage} the
choice constrains nothing. For quantitative researchers at the market makers the level-III median is below the
level-II one.

\begin{omfigure}
\begin{tikzpicture}
\begin{axis}[width=10.5cm,height=5.4cm,xmin=0.8,xmax=4.2,ymin=60,ymax=260,xtick={1,2,3,4},
  xticklabels={I,II,III,IV},tick label style={font=\footnotesize},
  xlabel={\small wage level: I entry, II qualified, III experienced, IV fully competent},
  ylabel={\small median offered base, \$ thousand},grid=major,grid style={gray!20},table/col sep=comma,
  unbounded coords=jump,legend style={legend cell align=left,font=\scriptsize,at={(0.5,-0.3)},anchor=north,
  /tikz/every even column/.append style={column sep=8pt}},legend columns=4]
\addplot[omDef,thick,mark=*] table[x=level,y=bank]{figdata/industry/14-pay-levels-by-role-firm-type-and-seniority/base_by_level.csv};
\addplot[omThm,thick,mark=square*] table[x=level,y=market_maker]{figdata/industry/14-pay-levels-by-role-firm-type-and-seniority/base_by_level.csv};
\addplot[black!70,thick,mark=triangle*] table[x=level,y=systematic_fund]{figdata/industry/14-pay-levels-by-role-firm-type-and-seniority/base_by_level.csv};
\addplot[gray,thick,dashed,mark=o,mark options={solid}] table[x=level,y=exchange]{figdata/industry/14-pay-levels-by-role-firm-type-and-seniority/base_by_level.csv};
\legend{bank,market maker,systematic fund,exchange}
\end{axis}
\end{tikzpicture}

\omcaption{Median offered base by prevailing-wage level, fiscal 2025, all roles. The level tracks pay at the banks and the
exchanges, and flattens above level~I at the market makers. The systematic funds' level-I cell is suppressed. Data: as
\cref{fig:in:pay-levels-by-role-firm-type-and-seniority:kind}.}\label{fig:in:pay-levels-by-role-firm-type-and-seniority:level}
\end{omfigure}

The occupational survey is the check from outside the visa system. In May 2025, software developers in the securities
and investments industry had a median wage of \$163\,290 (interquartile range \$132\,650 to \$203\,840), and in banks
(credit intermediation) \$136\,620; financial and investment analysts in securities \$124\,370. The filings' medians for
software engineers, \$175\,000 at the market makers, \$155\,700 at the banks and \$140\,400 at the exchanges, sit in the
same range: the sponsored positions of these employers are paid like the industry's, not below it.

\section{From base to total: what staff cost per head says about bonuses}

Base is the part the filings show; total pay is the part a candidate cares about. The only filed number that contains
the bonuses of a whole firm is staff cost per head, and the only filed number that describes a typical employee is
\omterm{def:in:pay-levels-by-role-firm-type-and-seniority:median}{median-employee pay}. Put against each other and against the visa filings, they bound the missing ratio.

\begin{method}[Implied ratio of total pay to base]\label{met:in:pay-levels-by-role-firm-type-and-seniority:ratio}
For a kind of employer, divide staff cost per head from filed accounts by the median offered base of the same kind in
the visa filings. The result bounds the ratio of total pay to base from above, because staff cost includes the
employer's social costs, benefits and deferred awards vesting, and it is a mean over a skewed distribution; and it mixes
two populations, all employees of one entity against sponsored positions of several.
\end{method}

\begin{example}[Three kinds of employer]\label{ex:in:pay-levels-by-role-firm-type-and-seniority:ratio}
A market maker's UK partnership spent between \$0.78 million and \$2.11 million per employee in 2023--2025 (chapter~11)
against a median offered base of \$175\,000 at the market makers: an implied ratio of 4.4 to 12.0. A London \omterm{def:in:quantitative-hedge-funds-and-systematic-managers:sys}{systematic
manager} spent \$4.11 million per employee in its year to January 2025 against the systematic funds' \$215\,000: 19.1. A US
bank's compensation and benefits were \$342\,000 to \$399\,000 per employee in 2023--2025 against the banks' \$154\,260:
2.2 to 2.6.
\end{example}

The ratio's range across kinds of firm is the chapter's central finding. At a bank, total pay is a small multiple of base
for the average employee, because the bank employs many people whose pay is mostly base; at the market maker and the
\omterm{def:in:quantitative-hedge-funds-and-systematic-managers:sys}{systematic manager}, whose staff are few and whose \omterm{def:in:revenue-and-profit-across-the-industry:perhead}{revenue per head} is high (chapter~12), the average is many times the
typical base. The caveats cut both ways: the entity's partners are outside its \omterm{def:in:reading-a-trading-firms-accounts:staff}{staff costs}, which understates the top,
and its average includes a few very large awards, which overstates the typical employee's.

\omterm{def:in:pay-levels-by-role-firm-type-and-seniority:median}{Median-employee pay} shows how much the average overstates the typical employee even where pay is moderate. For 2025 the
Intercontinental Exchange's median employee was paid \$112\,902 against staff cost per head of \$152\,834, 1.35 times as
much; Nasdaq's \$101\,885 against \$146\,142 (1.43); Morgan Stanley's \$136\,396 against \$352\,000 (2.58). A mean over a
bank's whole workforce is two and a half times its median employee's pay; in a firm whose pay is more concentrated, the
gap is larger.

\section{The tail: high-earner counts and risk-taker disclosures}

\begin{dated}{2026-04}{The best-paid in EU banks and investment firms}\label{dat:in:pay-levels-by-role-firm-type-and-seniority:tail}
\footnotesize
\begin{tabular}{@{}lrrr@{}}
\toprule
2024, EU & \omterm{def:in:pay-levels-by-role-firm-type-and-seniority:median}{high earners} & mean total pay (\euro) & variable over fixed\\
\midrule
banks, all areas & 2\,266 & 1\,830\,502 & 0.99\\
banks, investment banking & 904 & 1\,757\,914 & 1.23\\
investment firms, all areas & 288 & 2\,088\,931 & 3.59\\
investment firms, dealing on own account & 97 & 2\,094\,482 & 8.14\\
\bottomrule
\end{tabular}

\smallskip\noindent European Banking Authority, \omterm{def:in:pay-levels-by-role-firm-type-and-seniority:median}{high earners} 2024, Annexes II and III (published April 2026): 2\,554 \omterm{def:in:pay-levels-by-role-firm-type-and-seniority:median}{high
earners} in all (2\,343 in 2023). One UK bank's own report counts 717 of its staff paid \euro1 million or more in 2025
(chapter~11).
\end{dated}

The tail is where the kinds of firm differ most. Among EU \omterm{def:in:pay-levels-by-role-firm-type-and-seniority:median}{high earners}, those in banks' investment banking were paid
variable pay about equal to fixed on average, because the bonus cap still binds banks in the EU (chapter~15); those in
investment firms dealing on their own account, whom the cap does not cover, were paid variable pay 8.14 times fixed, and
the 97 of them averaged \euro2.09 million, about \$2.37 million at 2025's average rate. These are counts of people above a
threshold, so they say nothing about how many are near it, and they cover EU-authorised institutions only.

\begin{omfigure}
\begin{tikzpicture}
\begin{axis}[width=10.5cm,height=5.6cm,xmode=log,xmin=100,xmax=4000,ymin=0.4,ymax=6.6,ytick=data,
  yticklabels from table={figdata/industry/14-pay-levels-by-role-firm-type-and-seniority/slices.csv}{slice},
  yticklabel style={font=\scriptsize},xticklabel style={font=\footnotesize},log ticks with fixed point,
  xtick={100,200,500,1000,2000},xlabel={\small \$ thousand a year (log scale)},grid=major,grid style={gray!20},
  table/col sep=comma]
\addplot+[only marks,mark=*,mark size=2pt,color=omDef,error bars/.cd,x dir=plus,x explicit,
  error bar style={line width=2.5pt,omDef!50},error mark=none]
  table[x=lo,y=pos,x error expr=\thisrow{hi}-\thisrow{lo}]{figdata/industry/14-pay-levels-by-role-firm-type-and-seniority/slices.csv};
\end{axis}
\end{tikzpicture}

\omcaption{Four slices of pay on one log scale: offered base (interquartile range, visa filings, fiscal 2025), staff cost
per head from filed accounts (range over 2023--2025; one bank, one market maker's UK partnership), and the mean pay of
EU \omterm{def:in:pay-levels-by-role-firm-type-and-seniority:median}{high earners} in 2024 at 2025's average dollar rate. Each measures a different thing over a different population.
Data: \texttt{in\_paylevels} and the chapter's tables.}\label{fig:in:pay-levels-by-role-firm-type-and-seniority:slices}
\end{omfigure}

\section{Tutorial: the labour-condition pipeline}\label{tut:in:pay-levels-by-role-firm-type-and-seniority:tut}

\begin{tutorial}
\textbf{Goal.} Rebuild the chapter's ranges from the Department of Labor's files, and join them to the other slices.
\textbf{End state:} the range table,
\cref{fig:in:pay-levels-by-role-firm-type-and-seniority:kind,fig:in:pay-levels-by-role-firm-type-and-seniority:level,fig:in:pay-levels-by-role-firm-type-and-seniority:slices}.
\begin{tutsteps}
\item \textbf{Fetch.} The eight quarterly files (fiscal 2021 and 2025, about 1.2~GB) go to a scratch directory; they are
never committed.
\item \textbf{Stream.} \texttt{firm.paydata.read\_lca(path, keep)} yields one row at a time from the read-only workbook
(\cref{lst:in:pay-levels-by-role-firm-type-and-seniority:read}); the chapter's script
\texttt{in\_lca\_derive.py} caches the kept rows (kind, role, level, wage, nothing else) so the raw files are read once,
under a memory cap.
\omcode{code/firm/paydata/firm_paydata.py}{36}{50}{Stream a disclosure workbook row by row, keeping only the needed
columns.}\label{lst:in:pay-levels-by-role-firm-type-and-seniority:read}
\item \textbf{Classify and summarise.} \texttt{classify} with \texttt{data/industry/lca\_employers.csv},
\texttt{classify\_role}, \texttt{annualise} and \texttt{cell} (percentiles, bootstrap interval, suppression) write
\texttt{lca\_ranges.csv}, \texttt{lca\_titles.csv} and \texttt{lca\_counts.csv}.
\item \textbf{Join.} \texttt{in\_paylevels.implied\_ratios} and \texttt{median\_ratios} put the filings' \omterm{def:in:reading-a-trading-firms-accounts:staff}{staff costs} and
the proxies' medians against the ranges. The tests read only the committed tables and a synthetic workbook written at test
time.
\end{tutsteps}
\end{tutorial}

\textbf{What to change next.} Add fiscal 2022 to 2024 and follow each kind's median through the rate rises; split the
banks' software engineers by worksite state and compare with the state's occupational survey.

\section{Build: the pay-evidence reader}\label{bld:in:pay-levels-by-role-firm-type-and-seniority:paydata}

\begin{build}
\bfield{Purpose} Read the public pay sources into ranges and records that carry their caveats, for this chapter, the role
chapters (16--25) and chapter~30's choice.
\bfield{Interface} \texttt{firm.paydata}: \texttt{LCA\_COLS}; \texttt{read\_lca(path, keep)}; \texttt{annualise(amount,
unit)}; \texttt{Rule}, \texttt{load\_rules}, \texttt{classify(name, rules)}; \texttt{ROLE\_RULES},
\texttt{classify\_role(title, soc)}; \texttt{MIN\_CELL}; \texttt{cell(values, rng, n\_boot)};
\texttt{PayEvidence(source, measure, population, year, lo, hi, unit, caveat)}. openpyxl in read-only mode and NumPy.
\bfield{Rules} Stream, never load a whole file; read no contact or personal field; certified full-time applications
only; annualise the lower end of the offered range; suppress any cell under ten applications (the chapter's script adds
three employers); every source becomes a record with its measure, population and caveat.
\bfield{Acceptance tests} \texttt{code/firm/paydata/tests/}: a synthetic workbook is streamed and filtered without its
contact column; units annualise; employer and role rules classify and fall through to the SOC code; a nine-row cell is
suppressed and a hundred-row cell gives the right percentiles and an interval around its median.
\bfield{Stretch} The occupational survey by metropolitan area; PERM (permanent labour certification) filings, which carry
senior offers; a small-area estimate that borrows strength across kinds of employer.
\end{build}

\begin{omsources}
\item US Department of Labor, OFLC, LCA disclosure files for fiscal years 2021 and 2025, and the performance-data page;
20 CFR 655.731; OFLC Prevailing Wage Determination Policy Guidance (2009).
\item BLS, Standard Occupational Classification; O*NET 13-2099.01; OEWS May 2025 national industry-specific estimates
and technical notes.
\item 17 CFR 229.402(u); proxy statements for 2025 of the Intercontinental Exchange, Nasdaq, Morgan Stanley and Coinbase.
\item European Banking Authority, high earners 2024 and annexes; its legal notice.
\item The Goldman Sachs Group, Form 10-K for 2025; chapters 11--13 for the filed accounts and the Pillar~3 report.
\end{omsources}

\section{Exercises}

\begin{exercise}[$\star$]\label{exo:in:pay-levels-by-role-firm-type-and-seniority:1}
An application states a wage of \$95 an hour. What annual base does the chapter's rule record? And for \$16\,000 a month?
\end{exercise}

\begin{exercise}[$\star$]\label{exo:in:pay-levels-by-role-firm-type-and-seniority:2}
Using the table of medians, by how much does the systematic funds' median offered base for all roles exceed the
exchanges', in dollars and in per cent?
\end{exercise}

\begin{exercise}[$\star$]\label{exo:in:pay-levels-by-role-firm-type-and-seniority:3}
Compute the ratio of staff cost per head to \omterm{def:in:pay-levels-by-role-firm-type-and-seniority:median}{median-employee pay} for the Intercontinental Exchange in 2025, from \omterm{def:in:reading-a-trading-firms-accounts:staff}{staff
costs} of \$1\,963 million and 12\,844 employees.
\end{exercise}

\begin{exercise}[$\star\star$]\label{exo:in:pay-levels-by-role-firm-type-and-seniority:4}
Why can a bonus not be part of the wage in a labour condition application, and what does that make the filings measure?
\end{exercise}

\begin{exercise}[$\star\star$]\label{exo:in:pay-levels-by-role-firm-type-and-seniority:5}
The market makers' level-III median for quantitative researchers is below the level-II one. Give two explanations that
do not involve paying experienced people less.
\end{exercise}

\begin{exercise}[$\star\star$]\label{exo:in:pay-levels-by-role-firm-type-and-seniority:6}
Deflate the market makers' median offered base from fiscal 2021 to 2025 dollars with the consumer price index (114.33
in 2021, 135.83 in 2025) and compare with the fiscal 2025 median.
\end{exercise}

\begin{exercise}[$\star\star\star$]\label{exo:in:pay-levels-by-role-firm-type-and-seniority:7}
\emph{Coding.} Using \texttt{lca\_ranges.csv}, list every fiscal-2025 cell of the market makers that is suppressed, with
its number of applications and employers, and say which rule suppressed each.
\end{exercise}

\begin{exercise}[$\star\star\star$]\label{exo:in:pay-levels-by-role-firm-type-and-seniority:8}
\emph{Find the flaw.} ``The market makers' median base is \$175\,000 and their staff cost per head \$2.1 million, so a
new graduate there can expect total pay of about \$2 million.''
\end{exercise}

\section{Problem: Four Sources, One Number}

\begin{problem}[{Weekend problem --- four sources, one number}]\label{pb:in:pay-levels-by-role-firm-type-and-seniority:1}
A candidate choosing between a bank and a market maker wants one number: how much larger than base is total pay at each.
No source prints it. Build it.

\medskip\noindent\textbf{Part I --- The sources.}
\begin{enumerate}
\item Define a labour condition application and the \omterm{def:in:pay-levels-by-role-firm-type-and-seniority:lca}{prevailing wage}. What does the wage in the filing measure?
\item What does the occupational survey include and exclude from wages, and what does it top-code?
\item Define \omterm{def:in:pay-levels-by-role-firm-type-and-seniority:suppress}{top-coding} and \omterm{def:in:pay-levels-by-role-firm-type-and-seniority:suppress}{small-cell suppression}, and state the chapter's suppression rule.
\item Define \omterm{def:in:pay-levels-by-role-firm-type-and-seniority:median}{median-employee pay} and a \omterm{def:in:pay-levels-by-role-firm-type-and-seniority:median}{high earner}.
\item For each of the five sources, name what it measures and over whom.
\end{enumerate}

\medskip\noindent\textbf{Part II --- The ranges.}
\begin{enumerate}[resume]
\item State the method for building ranges from the filings.
\item Give the number of applications and employers by kind in fiscal 2025.
\item Give the median offered base by kind for all roles, quantitative researchers and software engineers.
\item How did the medians change from fiscal 2021 to 2025, nominally and after inflation?
\item What does the \omterm{def:in:pay-levels-by-role-firm-type-and-seniority:soc}{wage level} show at the banks and at the market makers, and why the difference?
\end{enumerate}

\medskip\noindent\textbf{Part III --- Base to total.}
\begin{enumerate}[resume]
\item State the method for the implied ratio of total pay to base and its biases.
\item Compute it for the market maker, the \omterm{def:in:quantitative-hedge-funds-and-systematic-managers:sys}{systematic manager} and the bank.
\item Compare staff cost per head with \omterm{def:in:pay-levels-by-role-firm-type-and-seniority:median}{median-employee pay} at three listed firms.
\item How does the occupational survey compare with the filings for software engineers?
\item What do the EU high-earner figures add, and what can they not say?
\end{enumerate}

\medskip\noindent\textbf{Part IV --- The verdict.}
\begin{enumerate}[resume]
\item State the \emph{named result}: the implied ratio of total pay to base by kind of employer, as a range with its caveats.
\item Which caveat most inflates the market maker's ratio for a typical employee?
\item Which slice would narrow the range most if it were public?
\item What should the candidate ask each employer to replace the ratio with facts?
\item In two sentences, answer the candidate.
\end{enumerate}
\end{problem}

\section{Interview questions}

\begin{interviewq}[$\star$ \iqroles{researcher, mle}]\label{iq:in:pay-levels-by-role-firm-type-and-seniority:1}
You have wages in five units (year, month, two weeks, week, hour). Write the conversion and say what you would do with a
range ``from \$150\,000 to \$250\,000''.
\end{interviewq}

\begin{interviewq}[$\star$ \iqroles{researcher}]\label{iq:in:pay-levels-by-role-firm-type-and-seniority:2}
Why report a bootstrap interval for a median rather than a standard error?
\end{interviewq}

\begin{interviewq}[$\star\star$ \iqroles{developer}]\label{iq:in:pay-levels-by-role-firm-type-and-seniority:3}
Design a job that reads a 300~MB spreadsheet in bounded memory and writes only aggregates. What can go wrong?
\end{interviewq}

\begin{interviewq}[$\star\star$ \iqroles{researcher, risk}]\label{iq:in:pay-levels-by-role-firm-type-and-seniority:4}
A cell has 12 observations from two employers. Should you publish its median? Why or why not?
\end{interviewq}

\begin{interviewq}[$\star\star$ \iqroles{bank, trader}]\label{iq:in:pay-levels-by-role-firm-type-and-seniority:5}
Why might variable pay be about equal to fixed pay for bank investment bankers in the EU, but eight times fixed pay at
investment firms dealing on own account?
\end{interviewq}

\begin{interviewq}[$\star\star\star$ \iqroles{researcher}]\label{iq:in:pay-levels-by-role-firm-type-and-seniority:6}
Visa filings cover only sponsored positions. How could selection bias their pay ranges, and how would you test for it?
\end{interviewq}
